\documentclass[11pt,letterpaper]{article}

\usepackage[margin=1in]{geometry}
\usepackage[T1]{fontenc}
\usepackage{lmodern}
\usepackage{microtype}
\usepackage{amsmath,amssymb}

\setlength{\parindent}{0pt}
\setlength{\parskip}{0.6em}

\title{Dielectric Insertion:\\Fields, Voltage, and Causality}
\author{}
\date{}

\begin{document}
\maketitle

Your cancellation argument is correct: for ideal infinite sheets at fixed free
charge, inserting the dielectric does not change the electric field in the air
gaps. However, \textbf{that field is not the entire electromagnetic field of the
capacitor--battery system}. The apparent contradiction arises when an
infinite-sheet \emph{electrostatic} result is used to rule out the fields needed
for a \emph{time-dependent circuit response}.

Saying ``the battery keeps the voltage constant'' gives the final equilibrium
condition, but does not explain how the battery responds causally.

\section{A voltage change does not require a field change at the plates}

Let $\sigma$ be the free surface-charge density on the positive plate, $d$ the
plate separation, $b$ the dielectric thickness, and $\kappa$ its relative
permittivity. In the ideal planar geometry, at fixed charge,
\[
    E_{\mathrm{air}}=\frac{\sigma}{\varepsilon_0},
    \qquad
    E_{\mathrm{dielectric}}=\frac{\sigma}{\kappa\varepsilon_0}.
\]
The dielectric's two infinite, uniformly charged bound surfaces cancel each
other's fields outside the slab, exactly as you argued. Nevertheless, the
positive plate's potential relative to the negative plate is
\begin{align*}
    V
    &=\int_{\text{positive plate}}^{\text{negative plate}}
        \mathbf{E}\cdot d\boldsymbol{\ell} \\
    &=E_{\mathrm{air}}(d-b)+E_{\mathrm{dielectric}}b \\
    &=\frac{\sigma}{\varepsilon_0}
        \left(d-b+\frac{b}{\kappa}\right).
\end{align*}
This is smaller than the original $\sigma d/\varepsilon_0$ because the field is
smaller \emph{along part of the path between the plates}.

There is no requirement for charges on the plates to experience a different
local force just because that integral changes. A potential difference is not
an additional local force that the charges must be ``notified'' about.

For the disconnected capacitor in the textbook, that is the whole story: the
voltage decreases, the free charge stays fixed, and no battery needs to respond.
A connected battery requires a more complete description.

\section{The battery responds to fields along its connection}

In an actual capacitor--battery circuit, inserting and polarizing the dielectric
involves moving bound charges and polarization currents. These, together with
the responding charges on the conductors, produce a changing electromagnetic
field.

The disturbance propagates through the apparatus at a finite speed. It changes
local electric fields along the leads and at the battery terminals, allowing
current to flow.

\textbf{The field driving current along a wire is its component along the wire.}
Your cancellation calculation concerns the field normal to the plates in the
ideal planar region. It does not establish that the field along the battery
connection remains unchanged.

Moreover, cancellation in an initial or final electrostatic state does not
establish cancellation throughout the transient. During that transient,
\[
    \oint\mathbf{E}\cdot d\boldsymbol{\ell}
    =-\frac{d\Phi_B}{dt},
\]
where $\Phi_B$ is the magnetic flux through the loop. The purely electrostatic
picture is therefore not the complete description of the response.

The battery does not remotely measure the dielectric and decide to compensate.
Its charges respond locally to electromagnetic and electrochemical forces.

\section{What changes if the plates are exactly infinite?}

Take the idealization completely seriously:
\begin{itemize}
    \item The plates and dielectric are exactly infinite in the directions
          parallel to the plates.
    \item Everything is translationally symmetric in those directions.
    \item The free surface-charge density is fixed.
\end{itemize}
Then the dielectric can indeed change the potential difference without changing
the external electrostatic field. But \textbf{that mathematical model contains
no localized battery connection whose response it can describe}.

A finite wire or contact, a hole through a plate, or a battery placed between the
plates breaks the symmetry. Alternatively, a connection going \emph{around}
truly infinite plates has no finite length. One cannot assume exact planar
symmetry while treating the geometry of a localized connection as irrelevant.

A closed-loop argument makes the missing information explicit. Compare the
\emph{fixed-charge electrostatic states before and after insertion} for a
\emph{finite} capacitor. Draw a closed path through the capacitor gap and back
around the outside. Write $\Delta\mathbf{E}$ for the difference between the two
static fields. Since both states are electrostatic,
\[
    0=\oint\Delta\mathbf{E}\cdot d\boldsymbol{\ell}
    =\underbrace{\int_{\mathrm{gap}}
        \Delta\mathbf{E}\cdot d\boldsymbol{\ell}}_{\text{nonzero}}
    +\int_{\mathrm{outside}}\Delta\mathbf{E}\cdot d\boldsymbol{\ell}.
\]
Consequently, \textbf{the field cannot remain unchanged everywhere along the
outside return path}. Nonuniform fields near and beyond the edges provide the
compensating contribution, even when the bulk air-gap field is essentially
unchanged.

For infinite plates, there is no finite outside return path of that kind. The
apparent puzzle comes from removing the edges and connection geometry while
still expecting the resulting model to explain communication with an external
battery.

The infinite-sheet approximation accurately describes the bulk electrostatics;
it is not a complete model of the circuit's causal response.

\section{Once the battery supplies charge, the air-gap field changes too}

The statement ``the air-gap field is unchanged'' is conditional on
\textbf{fixed $\sigma$}. For a battery maintaining the original voltage $V_0$,
the final equilibrium in the parallel-plate approximation instead satisfies
\[
    V_0=\frac{\sigma_{\mathrm{final}}}{\varepsilon_0}
        \left(d-b+\frac{b}{\kappa}\right).
\]
Thus,
\[
    \sigma_{\mathrm{final}}
    =\frac{\varepsilon_0 V_0}{d-b+b/\kappa},
    \qquad
    E_{\mathrm{air,final}}
    =\frac{V_0}{d-b+b/\kappa}.
\]
For $b>0$ and $\kappa>1$, this air-gap field is larger than the original $V_0/d$.
These equations describe the final state; the electromagnetic transient is how
a real circuit approaches it.

\section*{Bottom line}

You are right that causal influence must be electromagnetic. What fails is the
inference that an unchanged bulk air-gap field at fixed charge means no field
change can reach the battery.

\textbf{An unchanged field in one region does not imply unchanged fields
throughout the circuit, especially during insertion.} The ideal infinite-sheet
calculation omits the connection geometry and transient fields needed to answer
the communication question.

\end{document}
